\section{A fixed coherent weak space}\label{sec:fixed}
We prove Theorem~\ref{thm:fixed}. Write $\mathsf U\in\R^{g\times q}$ for an
orthonormal basis of $\im P$, so $\mathsf U\mathsf U^T=P$. One deterministic
choice is Gram--Schmidt on the first independent columns of $P$, in original
index order. If $u_\alpha\in\R^q$ is its $\alpha$th row, then
\begin{equation}\label{eq:frame}
 \sum_{\alpha=1}^g u_\alpha u_\alpha^T=I_q,\qquad
 \norm{u_\alpha}^2=P_{\alpha\alpha}\leq d:=\Lambda q/g.
\end{equation}
These coordinates serve the proof and channel selection. The observed fields
remain the original $z_{k,\alpha}$.

\subsection{The primary MSS input and an added count coordinate}
We use precisely the following classical theorem of
Marcus--Spielman--Srivastava \cite[Theorem 1.4, p.~3]{MSS}, in its version-pinned
arXiv statement.
\begin{theorem}[MSS finite-support theorem]\label{thm:MSS}
Let independent finitely supported random vectors $v_1,\ldots,v_m\in\C^N$
satisfy $\sum_i\Exp(v_iv_i^*)=I_N$ and
$\Exp\norm{v_i}^2\leq\varepsilon$ for every $i$, where $\varepsilon>0$.
With positive probability,
\[
 \norm{\sum_i v_iv_i^*}\leq(1+\sqrt\varepsilon)^2.
\]
\end{theorem}
We use this published theorem, not a proof or computational verification of it.
The simultaneous cardinality assertion needed next is not simply assumed
from a frame-partition corollary.

\begin{lemma}[Simultaneous Gram and count halving]\label{lem:halving}
Suppose $\sum_{i\in I}w_iw_i^T=I_q$, $m=|I|$, and
$\delta=\max_i\norm{w_i}^2+1/m$. There is a partition $I=S_1\sqcup S_2$
such that, for both $j=1,2$,
\begin{equation}\label{eq:half}
 (\tfrac12-e)I_q\leq\sum_{i\in S_j}w_iw_i^T\leq(\tfrac12+e)I_q,
 \qquad \tfrac12-e\leq |S_j|/m\leq\tfrac12+e,
 \quad e=\sqrt{2\delta}+\delta.
\end{equation}
All original rows, including zero rows, are counted.
\end{lemma}
\begin{proof}
Independently for every $i$, choose a uniform sign $\sigma_i\in\{-1,1\}$
and a uniform label $a_i\in\{1,2\}$, also independent within each pair.
Set $b_i=(w_i,\sigma_i/\sqrt m)\in\R^{q+1}$ and
$v_i=\sqrt2 e_{a_i}\otimes b_i\in\R^{2(q+1)}$.
The sign removes the off-diagonal covariance with the last coordinate:
\[
 \Exp(b_ib_i^T)=\diag(w_iw_i^T,1/m),\qquad
 \sum_i\Exp(v_iv_i^T)=I_{2(q+1)}.
\]
Each vector has four possible outcomes and squared norm
$2(\norm{w_i}^2+1/m)\leq2\delta$. Theorem~\ref{thm:MSS} therefore supplies
one common outcome with norm at most $(1+\sqrt{2\delta})^2$.
In each label block, the sum is $2\sum_{i:a_i=j}b_ib_i^T$.
Its leading principal block gives the upper Gram bound, and its last scalar
entry gives the upper count bound, both with
$(1+\sqrt{2\delta})^2/2=1/2+e$. Complementing the \emph{same} partition gives
both lower bounds. Signs do not alter the original Grams or cardinalities.
In particular a zero $w_i$ still has its count coordinate. The finite set of
feasible original-channel partitions has a lexicographically first member;
this specifies existence without claiming an efficient search.
\end{proof}

\subsection{Recursive whitening with a noncircular error bound}
Start with $I_0=\{1,\ldots,g\}$ and write
$G_t=\sum_{\alpha\in I_t}u_\alpha u_\alpha^T$, $m_t=|I_t|$.
When $G_t$ is positive, whiten \emph{only for selection}:
$w_\alpha=G_t^{-1/2}u_\alpha$. Apply Lemma~\ref{lem:halving}, choose a
deterministic child $I_{t+1}$, and set
\[
 \delta_t=\max_{\alpha\in I_t}\norm{G_t^{-1/2}u_\alpha}^2+1/m_t,
 \qquad\eta_t=2\sqrt{2\delta_t}+2\delta_t.
\]
Congruence by $G_t^{1/2}$ gives
\begin{equation}\label{eq:recurrence}
 \frac{1-\eta_t}{2}G_t\leq G_{t+1}\leq\frac{1+\eta_t}{2}G_t,
 \quad\frac{1-\eta_t}{2}m_t\leq m_{t+1}\leq\frac{1+\eta_t}{2}m_t.
\end{equation}

\begin{lemma}[Uniform conditioned prefixes]\label{lem:prefix}
Put $C_0=2^{24}$. If $4^k\leq g/(C_0\Lambda q)$, then $2k$ successive
halvings exist, and for every $0\leq t\leq2k$,
\begin{equation}\label{eq:prefix}
 \frac{15}{16}2^{-t}I_q\leq G_t\leq\frac{17}{16}2^{-t}I_q,
 \qquad\frac{15}{16}g2^{-t}\leq m_t\leq\frac{17}{16}g2^{-t}.
\end{equation}
\end{lemma}
\begin{proof}
Assume only the already established bounds at a current stage. Then
\[
 \delta_t\leq\frac{16}{15}2^t(d+1/g)
       \leq\frac{32}{15}2^t d,
\]
because $\Lambda q\geq1$. The current Gram is invertible, justifying that
stage's whitening. Over any justified prefix of $t<2k$, geometric sums give
\begin{equation}\label{eq:error}
 \sum_{t<2k}\eta_t\leq12\sqrt{d4^k}+5d4^k
 \leq\frac{12}{4096}+\frac5{2^{24}}<\frac1{256}.
\end{equation}
Here the displayed full sum denotes a majorant for every justified prefix,
not an assumption that future stages exist. Indeed
$\sum_{t<2k}2^{t/2}=(2^k-1)/(\sqrt2-1)$ and
\[
 \frac{2\sqrt2\sqrt{32/15}}{\sqrt2-1}<12,
 \qquad 2(32/15)<5.
\]
For example $\sqrt2<3/2$, $\sqrt{32/15}<3/2$ and
$(\sqrt2-1)^{-1}=\sqrt2+1<5/2$ already bound the first expression by
$45/4<12$. Also $d4^k\leq2^{-24}$; the last numerator in
\eqref{eq:error} is $49157<65536$ over denominator $2^{24}$.

For nonnegative errors with sum $S<1$,
$\prod(1-\eta_t)\geq1-S$ and
$\prod(1+\eta_t)\leq1/(1-S)$.
Start from $G_0=I_q,m_0=g$. The current loose bounds give the next error
bound; the resulting prefix products lie between $255/256$ and $256/255$,
strictly within $15/16$ and $17/16$. Applying \eqref{eq:recurrence} proves
\eqref{eq:prefix} at the next stage. This completes a prefix induction and
proves existence of all stages, rather than assuming their conditioning.
Every $\eta_t<1$, and both children have positive Gram and positive count.
The case $k=0$ contains no halving and is immediate.
\end{proof}

\subsection{Stopping, original groups and the full row count}
Set
\begin{equation}\label{eq:stop}
 M_* =\min\left\{\frac n{32},\sqrt{\frac g{C_0\Lambda q}}\right\}.
\end{equation}
If $M_*<1$, give \emph{all} channels singleton partitions. This priority
covers $n<32$ or $g<C_0\Lambda q$, including zero frame rows. It will be
counted and checked separately. Otherwise put $k=\lfloor\log_2M_*\rfloor$
and run exactly $2k$ halvings. For $0\leq j<k$ let
$J_j=I_{2j}\setminus I_{2j+2}$, and let $J_k=I_{2k}$ denote the one tail.
These disjoint sets exhaust all original channels. Their \emph{unwhitened}
Grams satisfy
\begin{equation}\label{eq:groupgram}
 G_{J_j}\geq\frac{43}{64}4^{-j}I_q\ (j<k),\qquad
 G_{J_k}\geq\frac{15}{16}4^{-k}I_q.
\end{equation}
The first constant is $(15/16)-(17/64)$, using a lower parent and upper
later subset in their Loewner difference. These are matrix bounds, not traces.

Use the actual spatial pieces of Lemma~\ref{lem:pieces}. Since
$2^k\leq n/32$, their parent lengths satisfy
\[
 L_j\leq\frac{33}{32}n2^{-j}\quad(0\leq j\leq k).
\]
Define $\lambda_j=5L_j^2/(8n^2)$, including the tail. Then
\begin{equation}\label{eq:lambda}
 \lambda_j\leq\frac{5445}{8192}4^{-j}
       <\frac{43}{64}4^{-j},
 \qquad \frac{43}{64}=\frac{5504}{8192}.
\end{equation}
Thus \eqref{eq:groupgram} implies $G_{J_j}\geq\lambda_jI_q$ for every
piece, also the tail. Give channels in $J_j$ the child partition
$\mathcal P_{j+1}$ when $j<k$, and channels in $J_k$ the singleton partition.
Every original channel receives exactly one partition. At $k=0$ all channels
belong to the singleton tail; $G_{J_0}=I_q\geq5I_q/8$ directly.

For $j<k$, $|J_j|\leq(17/16)g4^{-j}$. Hence, with the constant $A_D$ in
Theorem~\ref{thm:fixed},
\[
 N_{\rm coarse}\leq D2^D\frac{17}{16}g
              \sum_{j=0}^{k-1}2^{(D-2)j}
 \leq A_Dg2^{(D-2)k}\leq A_DgV/n^2.
\]
The sum is zero at $k=0$. Since $2^k>M_*/2$,
\[
 |J_k|\leq\frac{17}{16}g4^{-k}
 <\frac{17}{4}\max\{1024g/n^2,C_0\Lambda q\}
 \leq4352g/n^2+\frac{17}{4}C_0\Lambda q.
\]
Multiply by $DV$ for every singleton row to obtain
$N_{\rm out}\leq L_{D,\Lambda}(qV+gV/n^2)$.
In the priority branch the full count $DgV$ is at most
$1024DgV/n^2$ if $n<32$, or $DC_0\Lambda qV$ if $g<C_0\Lambda q$;
the same stated constant pays both. Equalities at either threshold use the
$k\geq0$ branch. No large-size inference is applied at an ineligible input.

\subsection{All-field coercivity with one strong-energy budget}
First take $z$ with componentwise mean zero. Put $f=\mathsf U^Tz$ and
$s=z-\mathsf U f$. These are $q$ weak forms and $g$ strong remainder forms,
respectively; $\mathsf U^Ts=0$. Because $\mathsf U$ is spatially constant,
\begin{equation}\label{eq:split}
 \norm z^2=\norm f^2+\norm s^2,\qquad
 Q_B(z)=a\En(f)+T,\qquad T=H\sum_\alpha\En(s_\alpha).
\end{equation}
For each actual group the parent-cell estimate gives
\[
 H\sum_{\alpha\in J_j}\En(s_\alpha)
 \geq\frac{4H}{L_j^2}\sum_{\alpha\in J_j}\norm{\Delta_js_\alpha}^2.
\]
The channel groups are disjoint: summing these inequalities spends $T$ once.
Since $\Delta_js_\alpha=\Delta_jz_\alpha-u_\alpha\cdot\Delta_j f$, use
\[
 \norm{p-v}^2\geq\frac9{10}\norm v^2-9\norm p^2.
\]
The difference is $10\norm p^2-2\ip p v+\norm v^2/10\geq0$.
Using the actual Grams and \eqref{eq:lambda} yields
\begin{align}\label{eq:strongbudget}
 T&\geq\frac9{10}\sum_{j=0}^k\frac{4H}{L_j^2}
       \ip{\Delta_j f}{G_{J_j}\Delta_j f}
       -36\sum_{j=0}^k\frac H{L_j^2}
                     \sum_{\alpha\in J_j}\norm{\Delta_jz_\alpha}^2\nonumber\\
 &\geq\frac94c\norm f^2-36cb\norm{\mathcal Rz}^2.
\end{align}
The positive pieces sum to $\norm f^2$ by mean zero and their orthogonality.
For the negative term, $\Delta_j$ is a subprojection of its channel's child
observation, or of $I$ at the tail. Also $L_j\geq n2^{-j}$ gives
$H/L_j^2\leq c4^j\leq c4^k$, and stopping implies $4^k\leq b$.
Disjoint channel groups therefore bound the negative sum by the displayed
single observation norm. No factor $q$ or extra copy of strong energy is used.

Independently Lemma~\ref{lem:fourier} gives $T\geq16c\norm s^2$.
Take the convex combination with weights $8/9$ and $1/9$ of this inequality
and \eqref{eq:strongbudget}. It gives
\[
 T+32cb\norm{\mathcal Rz}^2
 \geq2c\norm f^2+\frac{16}{9}c\norm s^2\geq c\norm z^2.
\]
Add $a\En(f)\geq0$. In the priority branch $\mathcal R=I$ and $K=c$
already suffice, so $32cb$ do also.

Finally every channel partition observes every global constant:
$\mathcal R\Pi_0=\Pi_0\mathcal R=\Pi_0$. Its constant and mean-zero images
are orthogonal, $Q_B$ ignores constants, and $32cb\geq c$ pays their full
norm. This proves the all-field statement and, by \eqref{eq:RC}, the precise
normalization in \eqref{eq:fixedupper}. Strong constants on zero weak-frame
rows are included, not assigned a nonexistent spectral gap.

\subsection{Actual-rank lower bound and the cancellation of weak dimension}
Let $C=\bigoplus_\alpha C_\alpha$ be any internally channel-separable linear
map on $\UU$. Let $S_\alpha$ be its row-space orthogonal projection
\emph{inside} $\UU_0$, and $S=\bigoplus_\alpha S_\alpha$. Then
$r=\rank(C|_\UU)=\sum_\alpha\rank S_\alpha$ and $\ker C=\ker S$.
The weak embedding $T_0:\UU_0^q\to\UU$, $T_0f=\mathsf U f$, is an isometry
of real domain dimension $qd_0$. Constant internal coefficients preserve the
actual coexact space. Lemma~\ref{lem:fourier} gives
\begin{equation}\label{eq:weakenergy}
 Q_B(T_0f)=a\En(f)\leq\frac c4\norm f^2.
\end{equation}
Taking a real orthonormal weak basis gives
\begin{equation}\label{eq:ranktrace}
 \tr(T_0^*ST_0)=\sum_\alpha P_{\alpha\alpha}\rank S_\alpha
       \leq\frac{\Lambda q}{g}r.
\end{equation}
Averaging over dimension $qd_0$ therefore supplies a \emph{unit} $f$ with
$t=\norm{ST_0f}^2\leq\Lambda r/(gd_0)$: the two $q$ factors cancel.
Set $w=T_0f$ and $u=(I-S)w$. Then $u\in\UU$, $Cu=0$ and
$\norm u^2=1-t$. Positive-semidefinite polarization gives
\[
 Q_B(u)\leq2Q_B(w)+2Q_B(Sw)\leq c/2+8DHt.
\]
The last step uses $Q_B(v)\leq4DH\norm v^2$ on actual fields.
At zero observation the restoring inequality, even with any finite
$\Phi$ satisfying $\Phi(0)=0$, implies
$c(1-t)\leq c/2+8Dcn^2t$. Thus
$t\geq[2(1+8Dn^2)]^{-1}$, proving the second bound in
\eqref{eq:ranklower} by \eqref{eq:ranktrace}.

Independently $CT_0$ must be injective, since a nonzero weak field in its
kernel would violate restoration by \eqref{eq:weakenergy}. Hence $r\geq qd_0$.
Since $d_0\geq(D-1)V/2$ and
$1+8Dn^2\leq(1+8D)n^2$, their maximum is at least
\[
 \frac{D-1}{8\Lambda(1+8D)}(qV+gV/n^2).
\]
Together with the actual canonical construction this proves the rank order.
Both a minimum on $\UU$ and a minimum requiring all-field restoration obey
these orders, because the lower is already on $\UU$ and the upper is all-field.
There is no output-dimension substitution in the lower argument.

\subsection{Least coefficient for the constructed fixed weights}
For a canonical map, $\mathcal R_\alpha|_{\UU_0}$ is a contraction into the
ambient one-channel field space; its squared singular values sum to at most
its actual rank. That rank equals $\rank(C_\alpha|_{\UU_0})$ by their equal
kernels in \eqref{eq:RC}, even if the projection image leaves $\UU_0$.
Consequently
\[
 \tr(T_0^*\mathcal R^*\mathcal RT_0)
 =\sum_\alpha P_{\alpha\alpha}
             \norm{\mathcal R_\alpha|_{\UU_0}}_{\rm HS}^2
 \leq\frac{\Lambda q}{g}r.
\]
Restoration on weak fields and \eqref{eq:weakenergy} require
$K\norm{\mathcal RT_0f}^2\geq(3c/4)\norm f^2$.
Sum a real orthonormal weak basis and cancel $q$ to obtain
\begin{equation}\label{eq:Ktrace}
 K\geq\frac{3cgd_0}{4\Lambda r}.
\end{equation}
For the constructed map $r\leq L_{D,\Lambda}(qV+gV/n^2)$, so
\[
 K\geq\frac{3(D-1)c}{8\Lambda L_{D,\Lambda}}
          \frac g{q+g/n^2}
 \geq\frac{3(D-1)c}{16\Lambda L_{D,\Lambda}}\min\{g/q,n^2\}.
\]
The last comparison follows by considering which denominator term is larger.
The sufficient all-field $32cb$ proves the opposite bound. A least coefficient
exists: the admissible set is a nonempty closed intersection of restoring
half-lines in $K$, bounded below by \eqref{eq:Ktrace}. This works on $\UU$
or all fields for the same map. These statements complete
Theorem~\ref{thm:fixed}.

\begin{remark}[Endpoints and the role of coherence]\label{rem:endpoints}
For $q=g$, $P=I$ and $M_*<1$; there is no strong complement. The singleton
map has actual rank $gd_0$ and sufficient $K=c$, with $b=1$. The singleton
priority also makes $n=2$, $n<32$ and $g<C_0\Lambda q$ explicit, rather than
an extrapolation of whitening. At $n=32$ or $g=C_0\Lambda q$, $k=0$ is
permitted and the one singleton tail is valid.
Without coherence a coordinate weak $q$-plane can use singleton weak channels
and one root on each strong channel, with actual rank $qd_0$ and $K=c$;
Lemma~\ref{lem:fourier} controls strong mean-zero forms and the roots pay
strong constants. But its diagonal entries equal one, so it is inadmissible
for fixed $\Lambda$ when $g>\Lambda q$. In the remaining range the extra
$gV/n^2$ term is controlled by a $\Lambda$-dependent multiple of $qV$.
This supplies no contrary coherent example.
\end{remark}
